% !TeX root = ../collection-solutions.tex

\titledquestion{Wealth taxes: why do the two papers disagree?}

\citet{guvenen2023use} find that a wealth tax of about $3\%$ is optimal and dominates capital income taxation. \citet{boar2023income-or-wealth} find that the optimal wealth tax is about \emph{zero} and that capital income taxation dominates — in a model built on the \emph{same} return-heterogeneity mechanism.

\begin{parts}
    \part[2] On what do the two papers \emph{agree}? State the shared mechanism.

    \begin{solutionorbox}[7cm]
        Both papers feature entrepreneurs with heterogeneous, persistent productivity who face collateral constraints, so returns to wealth differ across people and capital is misallocated. In both models a wealth tax improves the allocation of capital: it shifts the tax burden from productive to unproductive wealth holders (use it or lose it), and Boar--Midrigan confirm that in their economy, too, the wealth tax substantially reduces misallocation.

        Grading: 1 point for return heterogeneity from collateral-constrained entrepreneurs, 1 point for the shared use-it-or-lose-it efficiency channel.
    \end{solutionorbox}

    \part[4] Which modelling choices drive the opposite conclusions? Explain at least two, and how each tilts the ranking.

    \begin{solutionorbox}[12cm]
        \begin{itemize}
            \item (2 points) \emph{The size of the efficiency pie.} In Guvenen et al., all production runs through constrained entrepreneurs, so misallocation is large ($\approx 16\%$ by their Hsieh--Klenow measure) and reallocation gains are first order. Boar--Midrigan add an unconstrained corporate sector ($\approx 63\%$ of sales), give entrepreneurs a labor input with decreasing returns, and calibrate highly persistent ability — so misallocation losses are only $\approx 1.3\%$ of output, and the wealth tax has little efficiency gain left to harvest.
            \item (2 points) \emph{What redistribution the planner can do.} Boar--Midrigan let the government rebate revenue as lump-sum transfers, giving the planner a direct redistributive use for revenue raised from rich entrepreneurs' profits — this is what makes a high capital income tax attractive. In Guvenen et al.\ there is no lump-sum instrument; redistribution works indirectly through higher wages and a lower labor tax, a channel the wealth tax serves well.
            \item (alternative, also worth 2 points) \emph{Welfare criterion.} Boar--Midrigan compute welfare including the transition for the living population; Guvenen et al.'s baseline maximizes steady-state newborn welfare (their own transition extension weakens the case for the capital-income subsidy but preserves the wealth-tax gains).
        \end{itemize}
        Any two of the three, correctly explained, give full marks. The first (misallocation magnitudes) should be rewarded most generously — it is the heart of the disagreement.
    \end{solutionorbox}
\end{parts}
